\documentclass[12pt,letterpaper]{article}
\usepackage[utf8]{inputenc}
\usepackage[spanish]{babel}
\usepackage{geometry}
\geometry{top=2.5cm, bottom=2.5cm, left=2.5cm, right=2.5cm}
\usepackage{enumitem}
\usepackage{titlesec}

\titleformat{\section}{\large\bfseries}{\thesection}{1em}{}

\begin{document}

\begin{center}
	{\LARGE \textbf{Entrevista de Usuario}} \\
	\vspace{0.5cm}
	\textbf{Nombre:} Camila Rojas \\
	\textbf{Edad:} 34 años \\
	\textbf{Residencia:} El Tabo, Chile
\end{center}

\vspace{0.8cm}

\section*{Preguntas y Respuestas}

\begin{enumerate}[leftmargin=*]
	\item \textbf{¿Con qué frecuencia encuentra plásticos o basura en la arena al caminar por las playas de El Tabo?} \\
	      Casi siempre. Especialmente durante la temporada alta de turismo o los fines de semana largos, cuando las playas de El Tabo quedan llenas de envoltorios de comida y vasos de plástico después de la tarde.

	\item \textbf{¿Qué tipos de residuos específicos (microplásticos, colillas, vidrios, envases) son más visibles y constantes en la superficie de la playa?} \\
	      Se ve mucho plástico de un solo uso: botellas, tapitas y bombillas. También colillas cerca de las zonas de descanso, pero lo que más me preocupa son los pedazos pequeños de plástico de colores que se mezclan con la arena húmeda.

	\item \textbf{¿Ha notado que la basura permanece en la arena a pesar de las jornadas de limpieza municipales o comunitarias?} \\
	      Sí, totalmente. El camión recolector o los trabajadores limpian las zonas principales temprano por la mañana, pero solo se llevan lo grande. El plástico picado, las tapitas y los restos pequeños se van acumulando y quedan enterrados cuando sube la marea.

	\item \textbf{¿Existen zonas de la playa donde la acumulación de desechos sea tan crítica que se haya vuelto un rasgo permanente del paisaje?} \\
	      En las playas principales de El Tabo no tanto porque se limpian más seguido, pero si caminas hacia los sectores de roqueríos o las dunas de los alrededores, hay verdaderos microbasurales de botellas y latas que nadie retira.

	\item \textbf{¿Qué cambios físicos ha observado en la calidad de la arena (presencia de partículas extrañas o cambios de color) que confirmen una degradación del suelo?} \\
	      He notado que en algunas zonas de alta afluencia la arena se siente menos limpia, con un tono más grisáceo y una textura pastosa debido a los residuos orgánicos y aceites solares mezclados con partículas plásticas diminutas.

	\item \textbf{¿Ha tenido alguna experiencia personal con la contaminación en la costa?} \\
	      Sí, el verano pasado mi mascota ingirió un trozo de plástico que encontró enterrado en la arena mientras jugábamos cerca del agua en El Tabo. Tuvimos que llevarlo de urgencia al veterinario. Desde entonces soy mucho más consciente del peligro de los residuos pequeños.
\end{enumerate}

\newpage

\begin{center}
	{\LARGE \textbf{Mapa de Empatía}}
\end{center}

\vspace{0.8cm}

\section*{¿Qué ve?}
\begin{itemize}[leftmargin=*]
	\item \textbf{En su entorno:} Observa la acumulación constante de plásticos de un solo uso (botellas, tapitas, envoltorios) en la orilla del mar en El Tabo, sobre todo después de periodos de alta afluencia turística.
	\item \textbf{Acciones de otros:} Nota que el servicio municipal de limpieza se enfoca en las zonas centrales y en recolectar residuos de gran volumen, descuidando los plásticos fragmentados y las dunas.
	\item \textbf{Paisaje:} Ve que las rocas y sectores más apartados albergan acumulaciones de basura que no suelen ser cubiertas por los servicios básicos de aseo.
\end{itemize}

\section*{¿Qué dice y hace?}
\begin{itemize}[leftmargin=*]
	\item \textbf{Actitud:} Muestra una postura activa y de observación constante de su entorno natural inmediato como habitante de la comuna de El Tabo.
	\item \textbf{Conversación:} Comparte experiencias sobre el impacto directo de la contaminación en la vida cotidiana, relatando incidentes específicos de salud animal provocados por la ingesta de residuos.
	\item \textbf{Opinión pública:} Señala que existe una degradación progresiva de las playas (cambios en la textura y el aspecto de la arena) provocada por la acumulación histórica de microelementos y grasas.
\end{itemize}

\section*{¿Qué piensa y siente?}
\begin{itemize}[leftmargin=*]
	\item \textbf{Inquietudes:} Le preocupa la degradación silenciosa de la arena debido a partículas pequeñas y microplásticos que son difíciles de retirar de manera manual.
	\item \textbf{Sentimientos:} Experimenta frustración ante el comportamiento descuidado de los visitantes temporales y la falta de preservación del ecosistema local.
	\item \textbf{Percepción de eficacia:} Considera que los esfuerzos de mitigación actuales son superficiales, ya que no resuelven la acumulación de desechos de menor escala ni abarcan todas las áreas vulnerables de la costa.
\end{itemize}

\section*{¿Qué oye?}
\begin{itemize}[leftmargin=*]
	\item \textbf{Información del entorno:} Escucha quejas frecuentes de otros vecinos sobre el estado de la playa tras los fines de semana y está al tanto de las campañas locales de limpieza voluntaria que intentan cubrir los vacíos del servicio municipal.
\end{itemize}

\section*{Esfuerzos}
\begin{itemize}[leftmargin=*]
	\item \textbf{Frustraciones:} La persistencia de residuos plásticos fragmentados que resultan casi imposibles de limpiar con métodos tradicionales de recolección.
	\item \textbf{Miedos:} El riesgo continuo para la fauna local y las mascotas debido a la ingesta accidental de basura plástica oculta en la arena.
\end{itemize}

\section*{Resultados}
\begin{itemize}[leftmargin=*]
	\item \textbf{Necesidades:} Requiere de iniciativas de limpieza más minuciosas, con tecnologías o metodologías aptas para la recolección de microresiduos, además de mayor cobertura en zonas no turísticas.
	\item \textbf{Deseos:} Lograr un entorno costero verdaderamente preservado donde la biodiversidad local y las actividades recreativas en El Tabo no se vean comprometidas por la contaminación por plásticos.
\end{itemize}

\end{document}
